\chapter{系统非功能性与异常设计}

\section{鉴权与安全防护设计}
为了防范潜在的网络会话劫持与鉴权凭证伪造，LingoBridge 系统采用了基于数字签名的轻量级鉴权和安全防护机制。在鉴权生命周期中，所有的传输令牌均通过非对称加密算法进行哈希签名，其中包含了用户的全局唯一标识、已分配的角色列表以及绝对过期时间。客户端每次向受保护的接口发起调用时，请求拦截器会自动在请求头中注入该令牌。后端的网关中间件在接收到请求的第一时间会对其进行拦截，提取令牌后进行数字签名校验与时效比对。为了应对可能出现的令牌篡改，验证逻辑会强制匹配签名哈希，一旦检测到签名不符，中间件将拒绝执行后续的业务控制器，并直接返回未授权错误响应。在前端，全局响应拦截器一旦捕获到该错误响应，会立即触发安全清理流程，在内存和本地缓存中彻底清空当前的会话信息，并以原子性操作强制将页面重定向回登录页，从物理和逻辑两端防止了会话劫持风险。

\section{数据库写锁与并发争用应对策略}
由于平台采用轻量级的本地文件作为核心持久化载体，在高并发的多用户读写场景下，极易发生写争用锁死以及数据文件覆盖损毁的风险。为了化解这一工程痛点，数据层在内存缓存与落盘文件之间建立了一套基于排队机制的异步写队列与重试补偿策略。在数据访问组件中，所有针对内存数据的修改操作（如新增用户、更新学习进度等）首先在内存映射对象中被同步修改，以确保后续读操作的强一致性。接着，写操作会被作为一个待处理任务推送到全局唯一的串行写入队列中。该写入队列通过单线程轮询机制，每次仅允许一个写任务对底层的物理文件发起写入动作。为了应对偶尔出现的文件系统写锁争用异常，写入方法中深度嵌入了防抖退避重试算法。当物理文件由于操作系统锁占用导致写入失败时，引擎不会立即向前端抛出服务故障，而是自动将该次写入挂起，并在等待若干毫秒后再次尝试。若连续尝试多次均告失败，则将该任务标记为失效，并向调用上层抛出明确的内部服务器故障错误，由前端拦截器翻译为网络繁忙提示以引导用户稍后重试，从而彻底规避了高并发下文件数据损毁的可能性。

\section{课件文件传输异常与体积限制机制}
课件上传作为高数据吞吐量的业务，其传输过程中的稳定性和安全性直接关联到平台的整体稳定性。为了防止恶意的大体积文件上传导致服务器磁盘耗尽与内存溢出，系统在文件上传网关中设置了双重拦截防护堤。首先是在传输前端，文件选择器会在文件读入内存前，通过文件头元数据校验其媒体文件类型，严格限制仅允许演示文稿和便携式文档通过。同时，上传组件会实时读取文件的字节大小，一旦检测到文件体积超过五十兆字节，将立即终止上传，并在界面弹出带有警示标志的体积超限提示。其次在后端传输网关处，文件接收中间件同样会配置最大负载限制参数，若接收到的二进制流在传输中途超过此界限，中间件会立刻斩断套接字连接，并抛出实体过大的标准状态码。为了应对因网络波动造成的传输中断，文件网关提供了分片传输与网络异常断开重试的容错逻辑，确保在较差的无线网络环境中教师依然能够顺利将课件分发上屏。

\section{系统非功能性性能优化}
非功能性指标是评估在线中文学习平台是否卓越的核心维度，系统在设计阶段便针对网络延迟和渲染性能进行了全方位的优化。在文档渲染端，为了应对高清课件图像造成的内存占用过高以及缩放卡顿，系统引入了动态每英寸像素数适配策略。当用户缩放课件时，渲染引擎并不直接重新渲染全量文档，而是根据当前视口容器的实际缩放比例，计算出最优渲染率，仅在用户的缩放动作停止后触发一次局部的按需高清渲染，大大降低了图形处理器的功耗。在多媒体录音端，系统采用了经过高度压缩的音频编码标准，在保证发音测评所要求的声学特征不失真的前提下，使得录音文件体积缩减了百分之七十，显著降低了学生的上传带宽消耗。此外，数据层在内存中对高频访问的课程元数据和师生结对信息建立了二级缓存机制，使绝大部分读取操作的耗时控制在极低的毫秒级别以内，为全球中文学习者提供了顺滑、敏捷的交互体验。
